% APPENDIX 1. The systematic member-by-member specification. The Theory section carries the
% narrative (two closures, the ladder, the two structural differences) and points here for the
% formulas; this file carries the complete cycle with the three switches, so that a reader can
% reconstruct any member without reading the code.
% src: theory/macroir/docs/Macro_IR/macroir_derivation.tex (the canonical derivation: bilinear
%      tilde 740-768 and its closed form 880-941; vector tilde 1148-1165)
% src: legacy/qmodel.h:4498-4602 (the implemented cycle: the av=0 pre-propagation 4498-4512,
%      the two gSg forms 4525-4537, the variance-form switch 4543-4576, the two gains 4594-4602)
% NOT sourced from theory/macroir/docs/Macro_IR/MacroIR_tilde.md, which its own header marks as
% predating the corrected normalization.

\appendix
\setcounter{equation}{0}
\renewcommand{\theequation}{A\arabic{equation}}

\begin{appendixbox}
\section*{The interval update, member by member}

The Theory section gives the cycle for the boundary-conditioned member and names the two
places where the others differ. This appendix writes every member out in one template, so
that each is fixed by three switches on a single set of equations.

\subsection*{Objects computed once per interval}

All members share the quantities built from the rate matrix $\mathbf{Q}$ and the recording
filter for a window of length $\Delta$: the transition matrix $\mathbf{P}(\Delta)$; the
boundary-conditioned single-channel moments $\overline{\Gamma}_{i_0\to i_t}$ and
$\overline{V}_{i_0\to i_t}$; the start-conditioned mean conductance $\bar{\bm{\gamma}}_0$ of
Eq.~\ref{eq:gammabar}; the start-conditioned interval variance
$\bm{\sigma}^2_{\bar{\gamma}_0}$; and, for the members that ignore the acquisition
averaging, the instantaneous per-state conductance vector $\bm{\gamma}$. Two matrices
collect the boundary weights,
\begin{equation}
  G_{i_0 i_t} \;=\; \overline{\Gamma}_{i_0\to i_t}\,P_{i_0\to i_t}(\Delta),
  \qquad
  H_{i_0 i_t} \;=\; \overline{\Gamma}^2_{i_0\to i_t}\,P_{i_0\to i_t}(\Delta).
  \label{eq:GH}
\end{equation}

\subsection*{The tilde operator}

The tilde is a lift, a modulation and a collapse. An object carried on the $K$ states is
lifted to the $K^2$ boundary pairs $(i_0,i_t)$, multiplied there by the interval weights
$\overline{\Gamma}$, $\overline{V}$ and $\mathbf{P}(\Delta)$, and summed back down over the
pair index. Its two specializations are the ones the likelihood needs. The bilinear tilde is
the scalar of Eq.~\ref{eq:tilde}, which contracts the boundary covariance twice. The vector
tilde is the gain,
\begin{equation}
  \mathbf{g} \;=\; \widetilde{\bm{\gamma}^\top\bm{\Sigma}}
  \;=\; \mathbf{P}(\Delta)^\top\bigl(\bm{\Sigma}_0-\mathrm{diag}(\bm{\mu}_0)\bigr)\bar{\bm{\gamma}}_0
  \;+\; \mathbf{G}^\top\bm{\mu}_0,
  \label{eq:vector-tilde}
\end{equation}
which contracts it once and leaves a $K$-vector, the cross-covariance between the occupancy
at the end of the window and the interval current. Both are evaluated in $K\times K$
arithmetic; the $K^2\times K^2$ boundary covariance is never built.

\subsection*{One cycle, three switches}

Every member on the lattice runs the same four steps on each window: propagate the occupancy
Gaussian (Eqs.~\ref{eq:mu-prop} and~\ref{eq:sig-prop}), predict the observation
(Eqs.~\ref{eq:ypred} and~\ref{eq:pred-var}), update by the scalar Kalman correction
(Eqs.~\ref{eq:mu-post} and~\ref{eq:sig-post}), and either carry the posterior to the next
window or discard it. Three switches fix which member is being run.

\paragraph{Switch 1: the window, $\mathrm{av}\in\{0,1,2\}$.} It selects the conductance
object, and with it the second term of the predictive variance and the gain. At
$\mathrm{av}=0$ the observation is treated as an instantaneous sample taken at the centre of
the window: the occupancy is propagated there first, by $\mathbf{P}(\Delta/2)$, the conductance
is $\bm{\gamma}$, and, with $\bm{\Sigma}$ read in that mid-window frame,
\begin{equation}
  \widetilde{\bm{\gamma}^\top\bm{\Sigma}\bm{\gamma}} \;\to\; \bm{\gamma}^\top\bm{\Sigma}\bm{\gamma},
  \qquad
  \mathbf{g} \;\to\; \mathbf{P}(\Delta/2)^\top\bm{\Sigma}\,\bm{\gamma}.
  \label{eq:av0}
\end{equation}
At $\mathrm{av}=1$ the same two expressions hold with $\bar{\bm{\gamma}}_0$ in place of
$\bm{\gamma}$ and the full $\mathbf{P}(\Delta)$ in the gain, since the conductance is
conditioned on the state at the window's start and the whole window remains to be crossed;
the contraction is still an ordinary quadratic form on the $K$ states.
The trailing propagator in the gain is not decoration. It carries the update from the frame in
which the observation was formed to the frame the next window starts in, which is the end of
the current window, so it is the remaining half at $\mathrm{av}=0$ and the whole window at
$\mathrm{av}=1$. Omitting it, or using the full window where half is called for, leaves the
rank-one down-date in the wrong frame, and the symptom is specific and was observed: a spike in
the distortion-induced bias of the recursive instantaneous member at long intervals and large
channel counts.
% CORRECTED 2026-08-04. Eq. av0 read P(Delta) in the gain and named no frame for Sigma. Both are
% wrong for av=0: the implementation propagates to the window MIDPOINT and the trailing factor is
% P_half. Source, with the table that makes the three cases explicit and the historical failure it
% caused: docs/math/trust_coefficient.md:24-40 ("| 0 | X_mid via instantaneous g | g^T . Sigma_mid .
% P_half |", and "Without it the down-date is in the wrong frame - manifested historically as a
% Distortion-Induced-Bias spike in macro_R at long intervals / large Num_ch (fixed 2026-05-10)").
% Found by a replication audit, which noticed the appendix and that file disagree.
% CHECK BEFORE SUBMISSION: whether Delta/2 is exactly what the code uses, or whether P_half is
% defined against a different reference point (legacy/qmodel.h:4508-4515). At $\mathrm{av}=2$ the
conductance is conditioned on both endpoints and the two expressions become the tildes,
Eq.~\ref{eq:tilde} and Eq.~\ref{eq:vector-tilde}. The step from $\mathrm{av}=1$ to
$\mathrm{av}=2$ therefore does two things at once: it replaces $\bm{\Sigma}_0$ by
$\bm{\Sigma}_0-\mathrm{diag}(\bm{\mu}_0)$ and adds the boundary term ($\bm{\mu}_0^\top
\mathbf{H}\mathbf{1}$ in the variance, $\mathbf{G}^\top\bm{\mu}_0$ in the gain). The
subtracted diagonal is not a second approximation: the within-channel part it removes is
already inside the boundary term, and dropping the subtraction would count it twice.

\paragraph{Switch 2: recursion.} A recursive member carries
$(\bm{\mu}^{\mathrm{post}}, \bm{\Sigma}^{\mathrm{post}})$ into the next window. A
non-recursive one discards it and propagates the next window's predictive from the initial
condition, so its log-likelihood factorizes over windows. The switch changes no formula
above; it changes only what is fed to the next one.

\paragraph{Switch 3: the interval variance, total or residual.} It selects the third term of
Eq.~\ref{eq:pred-var}, which is present only when the acquisition averaging is modelled at
all ($\mathrm{av}\ge 1$):
\begin{equation}
  \bigl(\bm{\sigma}^2_{\bar{\gamma}_0}\bigr)_{i_0} =
  \begin{cases}
    \sum_{i_t} P_{i_0\to i_t}(\Delta)\bigl[\overline{V}_{i_0\to i_t}
      + \bigl(\overline{\Gamma}_{i_0\to i_t}-(\bar{\gamma}_0)_{i_0}\bigr)^2\bigr]
      & \text{total,}\\[4pt]
    \sum_{i_t} P_{i_0\to i_t}(\Delta)\,\overline{V}_{i_0\to i_t}
      & \text{residual.}
  \end{cases}
  \label{eq:variance-form}
\end{equation}
The two differ by the spread of the interval mean across end states, which is
Eq.~\ref{eq:mr-vr}. The switch is free only at $\mathrm{av}=1$. At $\mathrm{av}=2$ the
residual form is forced whatever the flag says, for the same reason the diagonal is
subtracted in switch~1: the between-endpoint spread is already carried by the boundary term,
so the total form would count it twice. % src: legacy/qmodel.h:4557 (averaging==2 || variance_form==residual selects the residual branch)

\subsection*{The members}

% Non-floating: a float inside appendixbox's mdframed can escape the frame.
\begin{center}
\begin{tabularx}{\linewidth}{@{}lccl>{\raggedright\arraybackslash}X@{}}
\toprule
Member & $\mathrm{av}$ & recursive & interval variance & what it is \\
\midrule
\texttt{NR}  & 0 & no  & --- & instantaneous samples, windows independent \\
\texttt{R}   & 0 & yes & --- & instantaneous samples, posterior carried forward \\
\texttt{INR} & 1 & no  & total & \texttt{NR} with the window average restored \\
\texttt{MR}  & 1 & yes & total & \texttt{R} with the window average restored \\
\texttt{VR}  & 1 & yes & residual & \texttt{MR} with switch~3 flipped, gain unchanged \\
\texttt{IR}  & 2 & yes & residual (forced) & \texttt{VR} with the boundary terms restored \\
\bottomrule
\end{tabularx}
\captionof{table}{The six members of the lattice as settings of the three switches. Reading down the
last column, each row differs from an earlier one in exactly one switch, which is what makes
the ladder measurable one step at a time. \texttt{IR} is MacroIR; \texttt{INR} is the
MacroINR of \citet{moffatt2025bayesian}. The data keys under which each is
dispatched are in Table~\ref{tab:family}.}
\label{tab:appendix-members}
\end{center}

One consequence of the two switches acting together is worth stating, because it is easy to
read the table as saying otherwise. \texttt{MR} and \texttt{IR} predict the \emph{same} total
variance. Expanding the total form of Eq.~\ref{eq:variance-form} and using
$\sum_{i_t}P_{i_0\to i_t}(\Delta)=1$ and $\sum_{i_t}P_{i_0\to i_t}(\Delta)\overline{\Gamma}_{i_0\to i_t}
=(\bar{\gamma}_0)_{i_0}$,
\begin{equation}
  \bm{\mu}_0^\top\bm{\sigma}^2_{\bar{\gamma}_0}\big|_{\text{total}}
  \;=\;
  \bm{\mu}_0^\top\bm{\sigma}^2_{\bar{\gamma}_0}\big|_{\text{residual}}
  \;+\; \bm{\mu}_0^\top\mathbf{H}\mathbf{1}
  \;-\; \bar{\bm{\gamma}}_0^\top\mathrm{diag}(\bm{\mu}_0)\,\bar{\bm{\gamma}}_0,
  \label{eq:mr-ir-identity}
\end{equation}
and the two extra pieces are exactly what switch~1 moves from the third term of
Eq.~\ref{eq:pred-var} into the second when $\mathrm{av}$ goes from $1$ to $2$. The
decompositions differ, the sum does not. What does change between the two members is the
gain, and that is the whole of the \texttt{MR}-to-\texttt{IR} difference.

Two of the six pairs isolate a single algebraic change and are the ones the Results use.
Stepping \texttt{MR} to \texttt{VR} flips switch~3 and touches nothing else, so it measures
what the variance form alone is worth. Stepping \texttt{VR} to \texttt{IR} restores the
boundary terms in both the variance and the gain and leaves the variance form where it
already was, so it measures what conditioning on the second endpoint alone is worth. The two
steps together are the \texttt{R}-to-\texttt{IR} gap.

\subsection*{Least squares, which is off the lattice}

Classical nonlinear least squares has no setting of the three switches. It propagates the
mean occupancy alone, $\bm{\mu}^{\mathrm{prop}} = \mathbf{P}(\Delta)^\top\bm{\mu}_0$, with no
covariance and no update, predicts $\bar{y}^{\mathrm{pred}}$ from Eq.~\ref{eq:ypred}, and
assigns one constant variance to every residual, fitted as the residual sum of squares over
the number of samples. Neither term of Eq.~\ref{eq:pred-var} that depends on the gating
survives, and there is no gain, so no occupancy object is ever conditioned on the data. It
is the anchor of the comparison rather than a member of the family.

\end{appendixbox}
